\documentclass[10pt,a4paper]{amsart}
\usepackage[margin=19mm]{geometry}
\usepackage{amsmath,amssymb,mathtools}
\usepackage[scr=rsfs,cal=euler]{mathalfa}
\usepackage{xfrac}
\usepackage[unicode,hidelinks,bookmarks=false]{hyperref}
\newcommand{\ie}{i.e.\ }
\newcommand{\cf}{cf.\ }
\newcommand{\resp}{respectively\ }
\newcommand{\Subsection}{\subsection}
\providecommand{\nobreakdash}{\nobreak\hbox{-}}
\setlength{\emergencystretch}{2em}
\multlinegap0pt
\usepackage{bm}
\usepackage{bbm}
\usepackage{dsfont}
\usepackage{xspace}
\usepackage{cancel}
\usepackage{mathtools}
\usepackage{amsthm}
\makeatletter
\@ifundefined{theorem}{\newtheorem{theorem}{Theorem}}{}
\@ifundefined{corollary}{\newtheorem{corollary}[theorem]{Corollary}}{}
\@ifundefined{lemma}{\newtheorem{lemma}[theorem]{Lemma}}{}
\@ifundefined{remark}{\newtheorem{remark}[theorem]{Remark}}{}
\makeatother
\newcommand{\CC}{{\mathbb C}}
\newcommand{\Per}{\mbox{Per} }
\newcommand{\SSS}{\mathbb S}
\newcommand{\cA}{{\mathcal A}}
\newcommand{\dd}{\mbox{\rm d}}
\newcommand{\mm}{\mathsf{m}}
\title[Wiener--Wintner points for topological dynamical systems]{Wiener--Wintner points for topological dynamical systems}
\author{Daniel Lenz, Nicolae Strungaru}
\date{Source: arXiv:2510.18056v1 (CC BY 4.0)}
\begin{document}
\maketitle

\noindent\emph{Source and licence.} Wiener--Wintner points for topological dynamical systems. Version arXiv:2510.18056v1 (CC BY 4.0), distributed under the Creative Commons Attribution 4.0 International licence, as recorded in the arXiv licence metadata for this exact version and shown on the abstract page https://arxiv.org/abs/2510.18056v1. Full rights notice: licenses/arxiv-2510.18056-CC-BY-4.0.txt.\par\smallskip

\noindent\textit{Editorial scope.} Counted unit: the Remark whose source label is \texttt{None} in the article (source0.tex lines 1020--1025, source label \texttt{None}), delivered together with the article's own complete original proof of it (source0.tex lines 1026--1200, one \texttt{proof} environment of 175 lines). No mathematics is written, repaired or re-derived: statement and proof are the article's own text line for line. Required context retained uncounted: lem-eigen-exi (lines 457--473); cor-van (lines 550--554). Every definition, display and label of those lines is retained unchanged. Omitted from this package and not redistributed: 1--456, 474--549, 555--1019, 1201--1593. Nothing inside a retained excerpt is omitted or altered: every retained body is delivered exactly as the article's lines are written, so no omission receipt of any kind accompanies this package. Citations: the retained excerpts carry no citation command at all, so no bibliography entry of the article is transcribed and no reference is redistributed. Reference adaptation: none was needed and none was made; every reference inside the delivered unit resolves inside it, so no unresolved reference or citation is produced. Printed slips kept literal: no printed word, symbol, hypothesis or number of the counted statement or of its proof was altered. Layout and class: the article's own class is \texttt{amsart} with packages including amsmath, amssymb, amscd, epsfig, amsxtra, pifont, mathabx, MnSymbol, accents, scalerel, stackengine, xcolor; the wrapper is the lane's pinned \texttt{amsart} editorial wrapper at 10pt on A4, which restates the theorem-like environments the retained text needs and adds the article's own macro definitions that the retained text needs (\texttt{\textbackslash CC}, \texttt{\textbackslash Per}, \texttt{\textbackslash SSS}, \texttt{\textbackslash cA}, \texttt{\textbackslash dd}, \texttt{\textbackslash mm}), with extra wrapper packages (bm, bbm, dsfont, xspace, cancel, mathtools). The delivered numbering is the unit's own. Assets: no retained span contains a figure environment, a graphics-inclusion command, a PDF-page inclusion, a table or a TikZ picture; no image file is copied into this package, the package declares no asset dependency and no image file is redistributed with it. Licence: version arXiv:2510.18056v1 of this article is distributed under the Creative Commons Attribution 4.0 International licence, as recorded in the arXiv licence metadata for this exact version and shown on the abstract page https://arxiv.org/abs/2510.18056v1. A full rights notice is delivered at licenses/arxiv-2510.18056-CC-BY-4.0.txt. Characters: every retained excerpt is pure ASCII, so the delivered engine, XeLaTeX with its default Latin Modern text font, reports no missing character.
\begin{lemma}[Main observation: forcing  eigenvalues] \label{lem-eigen-exi} Let $(X,G, \mm)$ be a measurable dynamical system, $\cA$ a F\o lner sequence and $G'$ a dense subgroup of $G$.

Let $E\subset \mathcal{L}^2 (X,\mm)$ be a $G'$-invariant subspace and let $x\in X$ be $(E,G')$-generic. Let $\bar{E}$ be the closure of $E$ in $L^2(X, \mm)$.

Let $(\chi_n)$ be a sequence in $\widehat{G}$  converging pointwise to $\chi \in \widehat{G}$ such that
$$ F(g):= \lim_n\frac{1}{|A_n|} \int_{A_n} g(t.x) \cdot \overline{\chi_n(t)} \dd t$$
exists for all $g$ in $E$. Then, there exists a unique  $f_\chi \in \bar{E}$  satisfying
\begin{itemize}
  \item{} $\| f_\chi \|_2 \leq 1$;
  \item{} $U_t f_\chi = \chi(t) f_\chi \,$ for all $t\in G$.
%  \item{} for all $g \in E$ we have
%  \[
% \langle g, f_\chi \rangle = a_\chi^{\mathcal{A}}(g_x) \qquad \mbox{ for all } g \in E.
%\]
\end{itemize}
Specifically, $f_\chi =0$ if $F(g) = 0$ for all $g\in E$ and  $f_\chi$ is an eigenfunction to $\chi$ otherwise.
\end{lemma}

\begin{corollary}[Vanishing of Fourier--Bohr coefficient of non-eigenvalues for generic points] \label{cor-van}
 Let $(X,G,m)$ be a measurable dynamical system and $\cA$ a F\o lner sequence and $G'$ a dense subgroup of $G$. Let $\chi\in \widehat{G}$ not be an eigenvalue.   Let $E$ be   $G'$-invariant with a countable core. Then,
$$\lim_n \frac{1}{|A_n|} \int_{A_n}  f(t.x) \cdot \overline{\chi(t)}  \dd t =0$$
holds for any $(E,G')$-generic point $x \in X$ and any $f \in E$. \qed
\end{corollary}

\begin{remark} The condition (iv) is equivalent to
\[
\| f- \sum_{j=1}^N c_j f_{\chi_j} \|_2= \|f_x - \sum_{j=1}^N c_j \chi_j  \|_2
\]
and the fact that the semi-norm on the RHS is actually the limit (and not just limsup).
\end{remark}

\begin{proof}
\noindent (i) $\Longrightarrow$ (ii) Since $x$ is generic, the existence of Fourier--Bohr coefficients for $\chi \notin \SSS$ follows from Corollary~\ref{cor-van}.


Let $\chi \in \SSS$ and let $e_\chi \in B$ be the eigenfunction for $\chi$. Let $f \in C(X)$. Then, by definition of Wiener--Wintner point, the Fourier--Bohr coefficient $a^\cA_{\chi}(f_x)$ exists and
\[
a^\cA_\chi(f_x)=\langle f, e_\chi\rangle \,.
\]
Let $g_\chi$ be any eigenfunction for $\chi$ be so that $\| g_\chi \|_2=1$. Since the eigenspace of $\chi$ is one dimensional by ergodicity, $g_\chi=c e_\chi$ holds for a suitable $c\neq 0$. Since both $g_\chi, e_\chi$ are normal, $|c|=1$.

Therefore,
\[
|\langle f, g_\chi \rangle|=|c| \cdot \langle f, e_\chi \rangle= |a^\cA_\chi(f_x)| \,.
\]



\smallskip

\noindent (ii) $\Longrightarrow$ (i)
Consider $\chi \in\widehat{G}$.  By Lemma \ref{lem-eigen-exi} we can restrict to the case that $\chi $ is an eigenvalue.

Pick some normalized eigenfunction $g_\chi$ and some $f \in C(X)$ so that $\langle f, g_\chi \rangle \neq 0$.

By Lemma~\ref{lem-eigen-exi} applied to $f$ and the core $E=C(X)$, there exists an eigenfunction $h_\chi$ so that
\[
a_\chi(g_x)= \langle g, f_\chi \rangle \qquad \mbox{ for all } g \in C(X) \,.
\]

%Define $F_\chi : C(X) \to \CC$ via
%\[
%F_\chi(f):=a_\chi^{\cA} (f_x) \,.
%\]
%Then, by (ii) we have
%\[
%\left| F_\chi(f) \right|^2= |\langle f, g_\chi\rangle|^2 \leq \| f\|_2^2
%\]
%Therefore, $F_\chi$ is a bounded operator on $C(X) \subseteq L^2(m)$. Since $C(X)$ is a dense subspace of $L^2(m)$, $F_\chi$ can be extended to a continuous functional on the Hilbert space $L^2(m)$. By Riesz' lemma, there exists some element %$f_\chi \in L^2(m)$ with $\| f_\chi \| \leq 1$ such that, for all
%$f \in L^2(m)$, we have
%\[
%a_\chi^{\cA} (f_x)=F_\chi(f) = \int_{X} f(y)\, \overline{ f_\chi(y)}\, \dd m(y) \,.
%\]
%By denseness of $C(X)$ in $L^2(m)$ we can find some $f \in C(X)$ with $\langle f,g_\chi\rangle\neq 0$ and from (ii) we find
%\[
%0 \neq |\langle f,g_\chi\rangle|^2 = |a_\chi^{\cA} (f_x)|^2=\left| \int_{X} f(\omega)\, \overline{ f_\chi(\omega)}\, \dd m(\omega) \right|^2 \ .
%\]
%This implies that $f_\chi$ is a non-trivial function. Moreover, for each $f \in C(X)$ a short computation gives
%$$
%\langle f, U_t f_\chi \rangle=\langle U_{-t} f, f_\chi \rangle = a^\cA_{\chi}((U_{-t} f)_x)=\overline{\chi(t)} a^\cA_{\chi}(f_x)= \overline{\chi(t)} \langle f,  f_\chi \rangle = \langle f, \chi(t)  f_\chi \rangle$$
%for all $f\in C(X)$. As $C(X)$ is dense in $L^2 (X,m)$  we get that $f_\chi$ is indeed an eigenfunction.

To finish this direction we show that the $(f_\chi)$ form a basis of $L^2 (X,\mm)_{\mathsf{pp}}$.

To do so it suffices to show that each eigenspace is one-dimensional. Assume that contrary and let $g,h$ be two orthogonal normalized  eigenfunctions to the same eigenvalue, say, $\chi$.

Then,
\[
|a_{\chi}^{\cA} (f_x)| = |\langle f, a g + b h\rangle| = |a \langle f,g\rangle + b \langle f,h\rangle|
\]
holds for all $f\in C(X)$ and all $a,b\in\CC$ with $|a|^2 + |b|^2 =1$.

Denoting by $z=  \langle f,g\rangle, w := \langle f,h\rangle$ we get that these two complex numbers satisfy
\[
|z|=|w|=|az+bw|
\]
for all $a,b \in \CC$ with $|a|^2+|b|^2=1$. It is straightforward to deduce that $z=w$.

This gives that
\[
\langle f,g\rangle = \langle f,h\rangle \,.
\]
Since this holds for all $f \in C(X)$, we get that $g=h$ which contradicts their orthonormality.



%This, however, is a contradiction as we could e.g. chose $f$ very close to $g$ and such force $\langle f, h\rangle$ to become arbitrarily small.

\smallskip


\noindent (i) $\Longrightarrow$ (iii) We first note  that the orbit $O_x=\{ t.x :t \in G\}$ is $G$-invariant. Therefore, by ergodicity $\mm(O_x) \in \{ 0,1 \}$. We split the problem into two cases:

{\it Case 1:} $\mm(O_x)=0$.

Then we can define
\[
h_\chi(y):=
\left\{
\begin{array}{cc}
  \chi(t) & \mbox{ if } y =t.x \in O_{x} \\
  f_\chi(y) &\mbox{ otherwise }
\end{array}
\right. \ .
\]
We show that this is well defined, and then it is immediate that $h_\chi$ is an eigenfunction satisfying the desired condition.
Pick one $f \in C(X)$ is such that $\langle f, f_\chi \rangle \neq 0$.
Assume that some $y \in O_{x}$ has two representations
\[
y=s.x = r.x \ .
\]
Then, the function $f_x$ is $s-r$ periodic. This gives
\[
\langle f, f_\chi \rangle = a_{\chi}^\cA(f_x)= a_{\chi}^\cA( ((U_{s-r} f)_x)= \overline{\chi(s-r)} \cdot a_{\chi}^\cA(f_x)= \overline{\chi(s-r)} \cdot a_{\chi}^\cA(f_x) \cdot \langle f, f_\chi \rangle \ .
\]
As $\langle f, f_\chi \rangle \neq 0$ we get $\chi(s)=\chi(r)$, completing the proof in this case.

\underline{Case 2:} $\mm(O_{x})=1$.

On $O_{x}$ the operation
\[
(t.x) \oplus (s.x) := (s+t).x
\]
defines a group structure ( and it is immediate that $(O_{x}, \oplus)$ is isomorphic as a topological group to $G/\Per(x)$).

It follows immediately that any $G$-invariant measure on this group is the Haar measure. Therefore, the ergodic measure is the unique probability Haar measure on $O_x$ and $O_x$ is a compact group.

Now, for each $f \in C(X)$ the function $f_x$ is $\Per(x)$-periodic and hence the Fourier--Bohr spectrum $\{ \chi : a_\chi(f_x) \neq 0 \}$ of $f_x$ is contained in
\[
(\Per(x))^0 \subseteq \widehat{G} \ .
\]
It follows immediately from here that for each eigenvalue $\chi$ the function we can define an eigenfunction
\[
h_\chi(y):=
\left\{
\begin{array}{cc}
  \chi(t) & \mbox{ if } y =\tau_t x \in O_{x} \\
  0 &\mbox{ otherwise }
\end{array}
\right. \ .
\]
continuous on the set $O_x$ of measure $0$.

Since $(O_x, G, \mm)$ is uniquely ergodic, and $h_\chi$ is continuous on $O_x$, we can apply the unique ergodic theorem on $O_x$ to get that for all $f \in C(X)$ and all eigenvalues $\chi$ we have
\[
\langle f, h_\chi \rangle =a_\chi^\cA(f_x) = \langle f, f_\chi \rangle \ .
\]
This implies that $\| f_\chi-h_\chi \|_2=0$ and hence $\| h_\chi \|_2=\|f_\chi\|_2=1$ and $h_\chi$ are indeed normalized.

\medskip
\noindent (iii) $\Longrightarrow$ (i) Let $h_\chi$, $\chi \in \SSS$,  be any such choice. Pick one $f \in C(X)$ such that
\[
\langle f, h_\chi \rangle \neq 0 \ .
\]
This implies that
\[
0\neq \lim_n \frac{1}{|A_n|} \int_{A_n} f(t.x) \cdot \overline{h_\chi(t.x)} \dd t
\]
and hence $h_\chi(x) \neq 0$.
Defining
\[
f_\chi:= \frac{1}{h_\chi(x)} h_\chi \ .
\]
we get a basis of normalized eigenfunctions associated to $x$ and $x$ is Wiener--Wintner point.


\noindent (i) $\Longrightarrow$ (iv) is immediate. Indeed
\begin{align*}
\| f- \sum_{j=1}^N c_j f_{\chi_j} \|^2_2 &= \int_{X} |f(y)|^2 \dd \mm(y)- \sum_{j=1}^N c_j \int_{X} \overline{ f(y)} f_{\chi_j}(y) \dd \mm(y)\\
&- \sum_{j=1}^N \bar{c_j} \overline{ \int_{X} \overline{ f(y)} f_{\chi_j}(y) \dd \mm(y)} + \sum_{j=1}^N |c_j|^2 \ .
\end{align*}
Since $f\in C(X)$ we can apply genericity for $x$ and $|f|^2$ and the definition of $x$ being Wiener--Wintner point for $\bar{f}$ to get
\begin{align*}
\| f- \sum_{j=1}^N c_j f_{\chi_j} \|^2_2&= \lim_n \frac{1}{|A_n|} \int_{A_n} \left(|f_x(t)|^2 - \sum_{j=1}^N c_j \overline{f_x(t)} \chi_j(t) \right. \\
&\left. -\sum_{j=1}^N \bar{c_j} f_x(t) \overline{\chi_j(t)} + \sum_{j=1}^N |c_j|^2 \right) \dd t\\
&= \lim_n  \frac{1}{|A_n|} \int_{A_n} \left|f_x(t) - \sum_{j=1}^N c_j \chi_j(t)  \right|^2 \dd t \ .
\end{align*}

\medskip (iv) $\Longrightarrow$ (i)
Note first that for each $f \in C(X), c \in \CC$ and $\chi$ we have
\begin{align*}
&\int_{X} |f(y)|^2   -   \overline{c} f(y) \overline{f_\chi(y)} - c \overline{f(y)} f_\chi(y) +|c|^2 \dd \mm(y) =\| f-  c f_{\chi} \|^2_2\\
&= \lim_n  \frac{1}{|A_n|} \int_{A_n} |f(t.x)|^2 -\overline{c} f(t.x)\overline{\chi(t)} - c \overline{f(t.x)} \chi(t) + |c|^2 \dd t \,.
\end{align*}
Applying this relation to $c=\pm 1, \pm i$ and taking a linear combination, via polarization we infer that $x$ is Wiener--Wintner point.
\end{proof}

\end{document}
